\section{Conclusion}
\label{sec:conclusion}

We presented a Hierarchical 5-Tier Framework for eye movement-based schizophrenia recognition that bridges clinical domain knowledge with deep representation learning. The framework progresses from raw fixation data through preprocessing, feature engineering of 135-dimensional trial-level expert vectors (including 60 category-contrast delta features), tabular GBDT baselines, two parallel deep architectures, and a meta-learner ensemble. Our proposed BiCA-HS model---combining a 4-layer Transformer encoder with a bidirectional cross-attention fusion of a 135-dimensional clinical expert stream---achieves an out-of-fold AUC of $0.9549 \pm 0.0039$, improving to $0.9590 \pm 0.0009$ with the Tier~5 logistic regression ensemble. Both results exceed the reported performance of the current state-of-the-art MSNet benchmark (reported test AUC\,=\,0.8854), with BiCA-HS maintaining clinical sensitivity of $87.5 \pm 4.7\%$ at the default operating threshold---a $40\%$ absolute improvement over the base ST-GNN.

Key findings include: (1) Tier~2B cross-category delta features encode stimulus-specific visual response differences directly linked to social cognition deficits, validated by SHAP analysis; (2) genuine non-degenerate cross-attention between the expert embedding and the pre-pooling fixation sequence (Direction~1) yields interpretable fixation-level clinical saliency maps; (3) the GNN-CEFAM architecture, despite Direction~2 degeneracy, converges stably and exceeds the reported MSNet SOTA, demonstrating that Direction~1's active attention provides sufficient representation capacity; and (4) BiCA-HS shows superior probability calibration, enabling threshold-free clinical deployment.

\textbf{Future work} includes: validating on independent multi-site datasets; extending graph representations to incorporate pupil dilation dynamics and face-gaze overlap statistics; replacing Direction~2 in CEFAM with an asymmetric attention over multi-head projections of the expert stream to eliminate the degeneracy; and cross-diagnostic experiments distinguishing schizophrenia from bipolar disorder and autism spectrum disorder.
